\clearpage
\section{TypeScript Integration Depth}

TypeScript quality is one of the clearest indicators of whether a UI library can serve a serious product team. Strong type systems improve discoverability, reduce regressions, and make complex component contracts easier to maintain over time.

\subsection{Generic component patterns in each library}
Generic patterns are essential when components must support multiple data shapes, layout variants, or display modes.

\begin{verbatim}
type TableProps<T> = {
  data: T[]
  renderRow: (row: T) => React.ReactNode
}

function Table<T>({ data, renderRow }: TableProps<T>) {
  return <>{data.map(renderRow)}</>
}
\end{verbatim}

\subsubsection{shadcn and Radix-oriented patterns}
Local source ownership usually gives teams the clearest type surfaces because component props are authored in-repo and can be constrained precisely.

\subsubsection{MUI patterns}
MUI commonly relies on strongly typed component props and theme augmentation. Large teams should keep wrapper types narrow so the API remains readable.

\subsubsection{Chakra UI patterns}
Chakra's prop model is ergonomic, but teams should be careful to avoid losing type precision when composing many helper wrappers.

\subsubsection{Ant Design patterns}
Ant Design type surfaces are broad and practical for enterprise data workflows. Teams often need local wrapper types to keep app code manageable.

\subsubsection{Tailwind and custom primitives}
Tailwind-oriented systems often pair with explicit prop interfaces and utility builders to recover type safety that would otherwise be hidden inside class strings.

\subsection{Prop type inheritance and extension}
Extending native element props is common, but teams should keep the inheritance surface intentional.

\begin{verbatim}
type ButtonProps = React.ButtonHTMLAttributes<HTMLButtonElement> & {
  tone?: "primary" | "secondary"
}
\end{verbatim}

\subsection{Discriminated union APIs}
Discriminated unions make complex component modes explicit and easier to validate.

\begin{verbatim}
type InputProps =
  | { kind: "text"; value: string }
  | { kind: "select"; options: string[]; value: string }
  | { kind: "date"; value: Date }
\end{verbatim}

\subsection{Theme type augmentation}
Theme augmentation is critical in MUI and similarly valuable in token-driven custom systems.

\begin{verbatim}
declare module "@mui/material/styles" {
  interface Palette {
    brand: Palette["primary"]
  }
  interface PaletteOptions {
    brand?: PaletteOptions["primary"]
  }
}
\end{verbatim}

\subsection{Event handler typing patterns}
Strong event typing reduces accidental mismatch between DOM state and component state.

\begin{verbatim}
function handleChange(e: React.ChangeEvent<HTMLInputElement>) {
  setValue(e.target.value)
}
\end{verbatim}

\subsection{Ref forwarding with TypeScript}
Ref forwarding is important for focus management, integration with form libraries, and composition with primitives.

\begin{verbatim}
const Input = React.forwardRef<HTMLInputElement, React.ComponentPropsWithoutRef<"input">>(
  function Input(props, ref) {
    return <input ref={ref} {...props} />
  }
)
\end{verbatim}

\subsection{Polymorphic component typing strategies}
Polymorphic `as` patterns are useful but should be wrapped carefully so the API remains understandable.

\begin{verbatim}
type AsProp<C extends React.ElementType> = {
  as?: C
}

type PolymorphicProps<C extends React.ElementType, Props = {}> = Props &
  AsProp<C> & Omit<React.ComponentPropsWithoutRef<C>, keyof Props | "as">
\end{verbatim}

\subsection{Using satisfies for theme configuration}
The `satisfies` operator helps keep literal values narrow while still enforcing shape.

\begin{verbatim}
const theme = {
  colors: { brand: "#1f2937" },
  radius: { sm: 4, md: 8 },
} satisfies ThemeConfig
\end{verbatim}

\subsection{Type-safe token systems}
Tokens should be modeled as unions or mapped types rather than loose strings.

\begin{verbatim}
type TokenName = "surface.primary" | "surface.muted" | "action.primary"
type TokenMap = Record<TokenName, string>
\end{verbatim}

\subsection{Autocomplete support across libraries}
Autocomplete quality depends on declaration quality, generic clarity, and the absence of too many indirection layers. Libraries with precise local wrappers often provide the strongest editor support.

\subsection{Common TypeScript errors and solutions}
\begin{table}[htbp]
  \centering
  \caption{Common TypeScript issues in UI libraries}
  \begin{tabularx}{\textwidth}{@{}p{0.28\textwidth}p{0.34\textwidth}Y@{}}
    \toprule
    \textbf{Issue} & \textbf{Typical cause} & \textbf{Typical solution} \\
    \midrule
    Excessively wide prop unions & Overextended wrapper abstractions & Split into discriminated variants or smaller wrapper layers. \\
    Theme augmentation errors & Missing module augmentation or mismatched declarations & Add explicit augmentation file and keep theme types centralized. \\
    Ref typing mismatch & Incorrect `forwardRef` generic order & Use explicit element and prop generics. \\
    Polymorphic component confusion & Too many `as` combinations & Restrict supported element types or publish helper aliases. \\
    Narrowing failures in conditional APIs & Missing discriminants & Introduce explicit `kind` or `variant` fields. \\
    \bottomrule
  \end{tabularx}
\end{table}

\subsection{Strict mode compatibility}
Strict mode exposes weak assumptions quickly. Libraries that tolerate loose `any` usage may appear easy to adopt, but they cost more when a team tries to harden types later.

\begin{itemize}[leftmargin=*, itemsep=0.35em]
  \item Keep prop definitions explicit and readable.
  \item Avoid implicit `children` assumptions where they are not intended.
  \item Use local wrapper types for especially broad external APIs.
  \item Prefer type-safe token keys over arbitrary string interpolation.
\end{itemize}

\subsection{Generic design patterns by library}
\begin{itemize}[leftmargin=*, itemsep=0.35em]
  \item shadcn/ui: strong fit for explicit local prop types and internal design tokens.
  \item Material UI: strong fit for theme augmentation and broad component wrappers.
  \item Chakra UI: strong fit for readable prop typing with careful wrapper discipline.
  \item Ant Design: strong fit for enterprise forms, tables, and workflow contracts.
  \item Radix UI: strong fit for primitive composition and clear interaction contracts.
  \item Tailwind UI: strongest when paired with a typed component layer built locally.
\end{itemize}

\begin{highlightbox}{TypeScript principle}
The best UI library for TypeScript is the one that lets teams express their product model precisely without turning the app code into a maze of type gymnastics.
\end{highlightbox}
